\documentclass[english, parskip=full]{scrartcl}
\usepackage[utf8]{inputenc}  % Kodierung der Datei
\usepackage[english]{babel}  % Sprache auf Deutsch 
\usepackage{color}
\usepackage{graphicx, gensymb}
\usepackage{amsfonts, cancel, amssymb, slashed}
\usepackage{amsmath, amsthm, array, esvect, esint}
\usepackage{booktabs, multirow}  % schönere Tabellen
\usepackage{caption}  % Anpassung der Captions
\usepackage{scrlayer-scrpage}    % Manipulation des Seitenstils
\pagestyle{scrheadings}  % Stil für die Seite setzen
\automark{section}  % Abschnittsnamen als Seitenbeschriftung 
\setheadsepline{.5pt}    % Kopzeile durch Linie abtrennen
\usepackage[hidelinks]{hyperref}  % Links und weitere PDF-Features
\usepackage{typearea, tikz, pgffor, verbatim}
\usetikzlibrary{calc, quotes}
\usepackage[cache=false, outputdir=build]{minted}
\usepackage{dsfont}


\definecolor{bg}{rgb}{0.9,0.9,0.9}
\newcommand\numberthis{\addtocounter{equation}{1}\tag{\theequation}}
\newcommand{\bra}{}
\newcommand{\kom}{}
\newcommand{\brak}{}
\newcommand{\braket}{}
\newcommand{\intx}{}
\newcommand{\intr}{}
\newcommand{\kla}{}
\newcommand{\klam}{}
\newcommand{\klammer}{}
\newcommand{\oper}{}
\newcommand{\tecxt}{}
\newcommand{\Bra}{}
\newcommand{\Ket}{}
\def\I{\text{i}}
\def\e{\text{e}}

\newcommand*{\titel}{10 Born approximation}
\newcommand*{\autor}{Joachim Schwardt}

\hypersetup{pdfauthor={\autor}, pdftitle={\titel}}

\subject{Quantum theory 2}
\title{\titel}
\author{\autor}
\date{\today}

\begin{document}
%\begin{titlepage}
\maketitle  % Titel setzen
%\tableofcontents  % Inhaltsverzeichnis setzen
%\end{titlepage}
%def convert_latex(string, dictionary=None):
%    if type(dictionary) == type(None):
%        dictionary = {0: {"old" : "align*", "new" : "align"}}
%    for i in range(len(dictionary)):
%        old_string = dictionary[i]["old"]
%        new_string = dictionary[i]["new"]
%        string = string.replace(old_string, new_string)
%    return string
\section{Questions}
\section{Scattering}
Let $V(x)=-V_0\e^{-r/r_0}$.
The scattering amplitude in first \textsc{Born} approximation with $q:=k-k'$ is
\begin{align}
f(\theta,\varphi)&=-\frac{m}{2\pi}\intr\int_{\mathbb{R}^{3}}V(x)\e^{-\I q\cdot x}\,\mathrm{d}^{3}x=\\
&=V_0m\intx\int_{0}^{\infty}\int_0^\pi\e^{-r/r_0}\e^{-\I qr\cos\theta}\,r^2\mathrm{d}r\mathrm{d}\theta=\\
&=V_0m\intx\int_{0}^{\infty}\e^{-r/r_0}\kla\left( {\e^{\I qr}-\e^{-\I qr}} \right)\frac{1}{-\I qr}\,r^2\mathrm{d}r=:\frac{V_0m}{-\I q}\kla\left( {I(q)-I(-q)} \right).
\end{align}
Here we defined 
\begin{align}
I(q)&:=\intx\int_{0}^{\infty}r\e^{-r/r_0}\e^{\I qr}\,\mathrm{d}r=\\
&=\klam\left[ r\frac{\e^{(-1/r_0+\I q)r}}{-1/r_0+\I q} \right]_0^\infty-\intx\int_{0}^{\infty}\frac{\e^{(-1/r_0+\I q)r}}{-1/r_0+\I q}\,\mathrm{d}r=\\
&=-\klam\left[ \frac{\e^{(-1/r_0+\I q)r}}{(-1/r_0+\I q)^2} \right]_0^\infty=\\
&=\frac{1}{(-1/r_0 +\I q)^2}=\frac{r_0^2}{1-q^2r_0^2-2\I qr_0}.
\end{align}
Inserting into the original expression yields
\begin{align}
f(\theta,\varphi)&=\frac{V_0m}{-\I q}\klam\left[ {\frac{r_0^2}{1-q^2r_0^2-2\I qr_0}-\frac{r_0^2}{1-q^2r_0^2+2\I qr_0}} \right]=\\
&=\frac{V_0mr_0^2}{-\I q}\frac{1-q^2r_0^2+2\I qr_0-1+q^2r_0^2+2\I qr_0}{(1-q^2r_0^2)^2+(2qr_0)^2}=\\
&=-\frac{4V_0mr_0^3}{(1+q^2r_0^2)^2}.
\end{align}
\section{Scattering by sphere}
Let $V(x)=-V_0\Theta (r_0-r)$.
The scattering amplitude in first \textsc{Born} approximation with $q:=k-k'$ is
\begin{align}
f(\theta,\varphi)&\simeq -\frac{m}{2\pi}\intr\int_{\mathbb{R}^{3}}V(x)\e^{-\I q\cdot x}\,\mathrm{d}^{3}x=V_0m\intx\int_{0}^{r_0}\intx\int_{0}^{\pi}\e^{-\I qr\cos\theta}\,r^2\sin\theta\mathrm{d}\theta\mathrm{d}r=\\
&=V_0m\intx\int_{0}^{r_0}\frac{1}{-\I qr}\kla\left( {\e^{\I qr}-\e^{-\I qr}} \right)\,r^2\mathrm{d}r=\\
&=-\frac{2V_0m}{q}\intx\int_{0}^{r_0}r\sin (qr)\,\mathrm{d}r=\\
&=-\frac{2V_0m}{q}\klam\left[ {-r\frac{\cos qr}{q}+\frac{\sin qr}{q^2}} \right]_0^{r_0}=\\
&=\frac{2V_0m}{q^3}\klam\left[ {qr_0\cos qr_0-\sin qr_0} \right].
\end{align}
Let $x:=qr_0$.
The differential cross section is
\begin{align}
\frac{\tecxt\text{d}\sigma}{\tecxt\text{d}\Omega}&=|f(\theta,\varphi)|^2=\underbrace{4V_0^2m^2r_0^6}_{=:\alpha}\klam\left[ \frac{x\cos x-\sin x}{x^3} \right]^2.
\end{align}
For $x\ll 1$ we can simplify this expression to
\begin{align}
\frac{\tecxt\text{d}\sigma}{\tecxt\text{d}\Omega}&= \alpha\klam\left[ {\frac{x\kla\left( {1-\frac{x^2}{2}+\frac{x^4}{24}} \right)-\kla\left( {x-\frac{x^3}{6}+\frac{x^5}{120}} \right)+\mathcal{O}(x^7)}{x^3}} \right]^2=\\
&= \alpha\klam\left[ {\frac{1}{3}-\frac{x^2}{30}+\mathcal{O}(x^4)} \right]^2.
\end{align}
This is similar to the \textsc{Yukawa} potential with $M\equiv\frac{1}{r_0}$ and 
\begin{align}
Y(x\ll 1)&\propto \frac{1}{1+x^2}=1-x^2+\mathcal{O}(x^4),
\end{align}
which leads to 
\begin{align}
\frac{\tecxt\text{d}\sigma}{\tecxt\text{d}\Omega}&\propto \klam\left[ {1-x^2} \right]^2.
\end{align}
\section{Green function for spherical waves}
Let $G(x):=\frac{-1}{4\pi r}\e^{-\I kr}$ with $r=|x|$. 
Then we can show that $G$ is the \textsc{Green}'s-function for the differential operator $D:=\nabla^2+k^2$, i.e. $DG=\delta$.
Knowing that $\nabla^2\frac{-1}{4\pi r}=\delta$ and $\nabla^2=\frac{1}{r}\partial_r^2r-\frac{\hat{L}^2}{r^2}$ yields
\begin{align}
\nabla^2G&=\nabla^2\frac{-1}{4\pi r}\sum_{n\ge 0}\frac{1}{n!}(-\I kr)^n=\\
&=\nabla^2\frac{-1}{4\pi r}+\nabla^2\frac{-1}{4\pi r}(-\I kr)+\frac{-1}{4\pi}\sum_{n\ge 2}\frac{(-\I k)^n}{n!}\kla\left( {\frac{1}{r}\partial_r^2r} \right) r^{n-1}=\\
&=\delta+\frac{1}{r}\partial_r^2\frac{\I k}{4\pi}+\frac{-1}{4\pi r}(-\I k)^2\sum_{n\ge 2}\frac{(-\I k)^{n-2}}{(n-2)!}r^{n-2}=\delta -k^2G.
\end{align}
Or, equivalently,
\begin{align}
(\nabla^2+k^2)G&=\delta.
\end{align}
\section{Fourier transformation}
The \textsc{Yukawa} potential is $Y(r)\propto \frac{\e^{-Mr}}{r}$.
Its \textsc{Fourier} transformation is
\begin{align}
\tilde{Y}(k)&=\mathcal{F}[Y](k)\propto \frac{1}{(2\pi)^3}\intr\int_{\mathbb{R}^{3}}\frac{\e^{-Mr}}{r}\e^{-\I k\cdot x}\,\mathrm{d}^{3}x=\\
&=\frac{1}{(2\pi)^2}\intx\int_{0}^{\infty}\intx\int_{0}^{\pi}\frac{\e^{-Mr}}{r}\e^{-\I kr\cos\theta}\,r^2\sin\theta\mathrm{d}\theta\mathrm{d}r=\\
&=\frac{1}{(2\pi)^2}\intx\int_{0}^{\infty}\e^{-Mr}\kla\left( \e^{\I kr}-\e^{-\I kr} \right)\frac{1}{-\I k}\,\mathrm{d}r=\\
&=\frac{\I}{4\pi^2k}\klam\left[ \frac{\e^{(-M+\I k)r}}{-M+\I k}-\frac{\e^{(-M-\I k)r}}{-M-\I k} \right]_0^\infty=\\
&=\frac{\I}{4\pi^2k}\klam\left[ {\frac{1}{M-\I k}-\frac{1}{M+\I k}} \right]=\\
&=\frac{\I}{4\pi^2k}\frac{M+\I k-M+\I k}{M^2+k^2}=\frac{1}{2\pi^2}\frac{-1}{M^2+k^2}.
\end{align}
\end{document}